% 原书第 117 页（PDF 第 130 页）
\BookPage{117}{130}
\BookHeading{《条款官话》初探}{《条款官话》初探}{赤岭守*}
长崎通事所学的一般是汉商使用的方言，与此不同的是，琉球国通事则一贯学习官话。较熟为人知的官话课本有《学官话》、《官话问答便语》、《白姓官话》《广应官话》四种。关于这四种官话课本的成书年代，濑户口律子氏在《琉球官話課本の成立時期に関する考証》中指出，《官话问答便语》编于1703年或1705年，《白姓官话》、《学官话》分别成书于1750年和1797年，《广应官话》则应是成书于1797至1820年间。《学官话》《官话问答便语》是琉球进贡官吏和勤学（琉球留学生）在闽与官吏、商人以及医生等进行问答的官话课本，《白姓官话》是基于1749年因遭风飘至琉球的白瑞临、瞿张顺等中国难民与琉球当地通事（翻译官）之间的问答内容编集而成的官话课本，三者之内容均探用问答形式。与上述三者不同的是，《广应官话》的内容非问答形式，而是按「天文」、「时令」、「地理」、「宫室」等30项分门别类编集而成。

\begin{BookQuote}
《学官话》、《官话问答便语》、《广应官话》是以学官话为主要目的编成的
\end{BookQuote}

课本，而同样作为课本的《白姓官话》却性质迥异。由于岛津氏入侵琉球后遂将琉球作为萨摩藩的附庸国纳入自己的统治下，并推行「隐蔽政策」，对中国隐瞒统治琉球的事实。在隐蔽政策展开过程中，萨摩岛津氏透过琉球的进贡贸易获得所需的中国物品。另一方面，为了维持王国体制，受统治于萨摩的事实不被中国发觉，首里王府也自发性地贯彻实施隐蔽政策。在推行隐蔽政策的过程中，王府最为担心的

\vfill
\begin{BookNotes}
\BookNote{*}{作者任职于琉球大学法文学部。}
\BookNote{1.}{口律子，〈琉球官話課本成立時期に関する考証〉，《官話問答語全訳—琉球官話研究》，榕树书林，2005年，172-173页。}
\end{BookNotes}

% 原书第 118 页（PDF 第 131 页）
\BookPage{118}{131}
\noindent 是琉球人飘到中国时和中国人飘到琉球时的应对问题。在琉球人飘至中国时的应对问题上，王府颁布了「旅行者心得之条」、「唐漂着船心得」、「御领国之船唐漂着之仪二付缔方」等条令，制定了全面否定日琉关系的应对方案。关于隐蔽工作的规定，「御船奉行」、「船改奉行」等诸官吏的公事帐上也均有明文记述。此外，飘至中国时应探取的隐瞒措施，《八重山岛地船上着公事帐并上国役人公事帐》(1859年）等地方的公事帐中亦有规定，隐蔽的指示不仅在琉球本岛，在各离岛也得到了贯彻执行²。中国人飘至琉球时，对负责飘风船只应对事务的「下知役」、「诘役」等官吏也同样作出了的指示和规定。《白姓官话》除具有官话教材的性质外，也是指示与中国飘风难民直接接触的通事如何隐瞒日琉关系的指南书。

本文所要介绍的《条款官话》，由文后「同治五年丙寅六月」的落款时间可知，其成书于1866年。若按官话教材的分类来看，则应与《白姓官话》一样，类属于对通事作出隐蔽指示，具有指南书性质的官话课本。文后还有「右は勅使樣御尋之節御返答之條官話組仕候事」等字句，显然是为应对册封正副使的种种提问而编的。「册封」按其形式可分为「颁封」和「领封」两种，中国派遣使节前往册封的形式称「颁封」，反之，被册封国遣使中国祇领封王诏敕的形式则为「领封」。琉球一贯接受中国遣使册封的「颁封」仪式，其册封是以中国皇帝至高无上的皇权为背景的，通过册封以确认琉球国王位继承的正当性。经过中国皇帝认可其统治权的政治过程，从而产生了政治统治上不可或缺的「权威」之向心统合力。换而言之，册封是表示王权得到中国皇帝承认的一种政治礼仪，但对琉球而言其重要意义还在于，与中国间的宗藩关系体系中琉球作为「藩属国」的地位可以得到认可，此关乎琉球王国的存续。有清一代，中国遣使往封琉球八次，使节一行人数多达四、五百人，逗留期间也长达四、五个月。在使节一行逗留期间，首里王府为了防止日琉关系败露进行了种种隐蔽工作。关于这些隐蔽工作，喜舍场朝贤在《琉球见闻录》中叙述如下：

\begin{BookQuote}
我邦和日本的关系一旦败露势必有碍进贡，故对中国隐瞒。册封钦差（敕使谓
\end{BookQuote}

\vfill
\begin{BookNotes}
\BookNote{2.}{渡辺美季，〈清に対する琉日関係の隱蔽と漂流問題〉，《史学雑誌》第114篇第11号，2005年，1807页。}
\end{BookNotes}

% 原书第 119 页（PDF 第 132 页）
\BookPage{119}{132}
\begin{BookQuote}
之钦差）来琉之时，在琉的萨摩官吏及商人全部移居浦添间切城间村，原停靠于泊港的萨摩船只亦移泊今归仁间切运天港，以避中国人耳目。平常国中流通的宽永钱悉数收匿，取出官库中备用的鸠目钱，在中国人逗留期间暂时流通于市。官衙寺院里的挂轴钟铭以及碑文等，凡有涉及日本年号、人名或者与日本有实质关系之物，尽数收匿。国中所用的物品器具，若是日本产之物则称是产于度佳刺岛。昔日，日中间皆锁港，互不通往来，不遣聘问使，故不知我邦内属于萨摩之情。3
\end{BookQuote}

为了贯彻隐蔽政策，王府颁布了〈冠船渡来二付取締方書渡候覚〉，下令册封使节逗留期间禁止首里、那霸的周边地方悬挂通告，与日本年号、日本姓名甚至日本书籍等有关的一切事务一概不许让使节一行看到。并且在若狭町、泊高桥和安里桥（崇元寺桥）三处设关卡，监视往来行人，禁止一般民众进出使节一行的下杨馆舍；不论日本产或是琉球土产之物，均禁止私自出售册封使团人员。此外，对逗留期间的中国人有可能提出的种种问题进行设想，颁布了〈唐人江万問答並晴樣の条〉，对负责应对的诸官吏以及相关人员作出彻底隐蔽的指示。

文后所记的「同治五年」乃琉球末代国王尚泰（1834-1901）受封之年，当时的册封正副使分别是赵新和于光甲4。因尚育王去世，年仅六岁的尚泰于道光二十八年（1848）即位成为了第二尚氏王统第十九代国王，但由于中国国内太平天国之乱爆发，册封延期直至叛乱被镇压之后才进行。在此期间琉球经历了一段多事之秋。在外交方面，英国传教士伯德令以及法国传教士伙尔加助、伯多禄强行逗留琉球；华工反乱事件的发生（罗伯特·包恩号事件）；英法船只来琉强行要求通商；逼迫日本开国的美国总督贝理率舰队来琉；琉美修好条约、琉法修好条约、琉兰修好条约的签订；岛津齐彬主导实施的以琉球为仲介的对法贸易计划等。在内政上，岛津

\vfill
\begin{BookNotes}
\BookNote{3.}{喜舍场朝贤，〈琉球见闻录〉，至言社，1977年，89页。}
\BookNote{4.}{册封正使赵新，福建省侯官县人，字又铭，号古彞，咸丰二年进士（三甲一二八名合格者中第一百零四名），受派遣时官任詹事府右春坊右赞善。册封副使于光甲天津府沧州人，字申卿，咸丰六年进士（二甲一百名合格者中第二十七名），受派遣时官任内阁中书舍人。（曾焕棋，《清代使琉球冊封使の研究》，榕树书林，2005年，311页、320-321页。夫马进，《使琉球錄解題及び研究》，榕树书林，1999年，135-137页。)}
\end{BookNotes}

% 原书第 120 页（PDF 第 133 页）
\BookPage{120}{133}
\noindent 齐彬死后围绕三司官人事案而引发得疑狱事件（牡志·恩河事件）、饥荒以及农村经济慢性疲敝现象的出现、一八六零年代发生的由铜钱、铁钱兑换汇率的变动而引发的琉球国财政危机和经济混乱等。

在尚泰的册封问题上，隐蔽工作已不单偏限于原来的日琉关系了。全力向中国展现从属意识的王府担心，一系列未经皇帝奏准的外交处理以及内政上的不恰当措施会有损宗主国皇帝的威严，有可能对册封甚至王国体制造成影响。因此，王府认为除对日琉关系进行隐瞒外，还有必要对敕使（册封正副使）可能问及之有关琉球外交、内政的种种问题进行设想，并将提问的标准回答向通事作出指示。继1854年日美和亲条约签订后，1858年又签订了日美修好通商条约，紧接着与荷兰、英国、法国、俄罗斯等国的条约也相继签署。随着这些条约的签订，在日外国人逐渐增加，外国船只也频繁出现于琉球海域。因此，琉球所探取的隐蔽政策和制度由海外泄露到中国的可能性也是存在的。

\begin{BookQuote}
《条款官话》就是在这种历史背景下编成的。赵新、于光甲等册封使一行
\end{BookQuote}

于同治五年（1866）6月22日抵琉，同年11月10日返国，在琉球的逗留时间长达138天，《条款官话》应是赶在册封使到来之前成书的。与其说《条款官话》是一本学官话的课本，倒不如说是一部探用问答形式编集而成的指南书，意在对通事作出隐蔽的应对指示，其作为指南书的浓厚色彩比起《白姓官话》而言，完全是有过之而无不及。现将其全文抄录如下：

\begin{BookQuote}
问：番洋船只到琉球的时节用得官音么？用得番洋说话么？5答：我们敝国原来不通西洋的话，只用官音替他对讲，若有不通官话的人，就
\end{BookQuote}

\vfill
\begin{BookNotes}
\BookNote{5.}{1842年中英签订南京条约，香港被纳入英国殖民统治下，上海、宁波、福州、厦门等五港也被迫开放通商。之后，西洋各列强亦纷纷逼迫中国签订各种不平等条约，来航中国的外国船只骤增。琉球也因海域探测之需要以及作为开放日本门户之踏板而受瞩目，英、法、美等外国的船只相继而来。1844年4月法国的阿鲁库美奴号航抵琉球，要求首里王府通商贸易以及传教，并不顾王府的反对强行把传教士伙尔加助等人留在琉球。1846年4月英国船史塔林号来航，传教士伯德令一家亦滞留琉球不归。同年6月，法国船库烈欧巴托奴号来航，进行通商、传教之条约交涉。此外，1853年5月又有贝理舰队来航，强行到访首里城。鸦片战争后，外国船只来航次数多达60次以上。在双方交涉过程中，王府也完全表现出对中国的属国意识，彻底隐蔽日琉关系。进行外交时，也仿照中国的官制，虚设总理官、布政官和地方官，取代摄政、三司官和那霸里主之职。}
\end{BookNotes}

% 原书第 121 页（PDF 第 134 页）
\BookPage{121}{134}
是用了手样，略略通得。问：前年番洋船只来到琉球，他们请见你的官府，那时节国相法司官都是相见

他么？答：呵，就是这样的。问：我听见说琉球的旗号递给番洋的人，这个事是真的么？答：这是有的。那番洋头目说，把你旗号送给我们作证，我把这旗号通知西洋

船人，看见这号船只遇着灾难，叫他救灾恤患不要推辞。所以写着巴字旗号给他作证。问：那个旗号给他以后，救灾有没有？答：有的。咸丰五年贵国的人飘到敝国，解送福建那船回国的时节半洋遇着飓

风，把船打坏，通船的人也有靠着船板也有靠着杠棋随风漂流，遇着西洋船只，被他救得活命。此年，又有敝国船只驶到太平山的时候遇着暴风，正是万死一生，幸是遇着西洋船只，也被他救的。咸丰六年敝国进贡的船，五虎门开洋遇着灾难，那船上数十人等移登西洋船只才得活命。问：这样灾难他替你们救起，这样恩德怎么样报得呢？答：那船只到国得时节送给礼物道谢。问：我听见说你们琉球做得文章递给英国军机大臣，又有做了文书屡屡送给西

洋人，这个事是真的么？答：敝国法度原来没有做文书送给外国人，只有英国官员做了文书访问接待留

国伯德令的事，没奈何做文报答。又有西洋船主投文告事，就是做文答应。也有敝国要撤回英佛得人做文投告，其余的事情并没有做文投告。问：我听见讲那伯德令留在琉球的时节，把那耶稣教必要教导球人，球人不肯

\vfill
\begin{BookNotes}
\BookNote{6.}{伯德令（1811-1870）是匈牙利出生的犹太人，也是由英国海军军人琉球传道会所派遣之新教传教士。1846年4月30日偕同其家人在那霸强行登陆，暂住于波之上护国寺，开始进行传教。但由于传教活动受到王府的多方阻碍，九年后的1854年7月17日伯德令无奈被迫放弃传教，搭乘贝理舰队的船只离开琉球。伯德令虽然传教失败，但却将预防天花的种痘法传入琉球，王府甚至还要求伯德令提供有关天花疫苗以及种痘法的册子抄本。此外，伯德令还在通事的协助下，将《圣书》翻译成琉球语在香港发行。（参考照屋善彦着、山口荣铁·新川右好译，《英宣教医べッテル八イム一琉球伝道の九年間》，人文书院，2004年)}
\end{BookNotes}

% 原书第 122 页（PDF 第 135 页）
\BookPage{122}{135}
受教，这话是真的么？答：那个伯德令天天把那耶稣教要教球人细细解劝。我们敝国往年以来学习四

书五经，多蒙天朝的教化，可以那耶稣教人心不服，不肯领教。问：我听见说那伯德令把那医生的事业、天文地理要教球人，球人不肯受教。

我想起来那耶稣教不是圣贤的道理，推辞更好。那医生的事业、天文地理都是要紧的事，学习更好。为何辞掉？答：敝国医生事业原来到中国去学，那个操舟楫、看天气的方法都是学得备

用，所以推辞。问：我听见说那伯德令回国时候你们琉球送他贰千五百两银子。为什么那多银

子送给他呢？答：那伯德令说我在本国时候把我洋银借给我的朋友，现今没有还，我若是空

身回国，家里伙食也是艰难的，等他银子还我然后回国。那伯德令若是好久留在敝国大有不便，要他快快回国，所以把他交给敝国的物价银子都是仍旧算明还他，他领了回国。问：我听见讲，你们琉球把那日本的说话教了番洋的人，这是真的么？答：我们敝国原没有替日本交通，所以没有人知道日本的话，只有那佛英的人

告禀说，我要学经书，替我教导。这个事不得推辞，就把经书都用球话教的，并没有教了日本的话。问：那佛国的人留在琉球以后，他天天做什么事情？又是告禀官府衙门，这事

有没有？答：他天天看见洋书或者学习讲话，不过这两样的事情。他要教天主教就有禀

上官府，或者交待的事情不合他的心思也有告禀的。这时候看他行状留心办事没有过失。问：我听见说日本的人改换琉球的粧扮假说大岛的人，一同琉球通事到的那佛

人房子，替他谈谈不觉讲出日本的话，是真的么？答：这不是真话。替那佛人交接来往这是违禁的。只许通事交接，不许别人往


% 原书第 123 页（PDF 第 136 页）
\BookPage{123}{136}
来。那大岛的人到不得他家。问：我听见说，番洋的人飘到这里，你们琉球报知留在长崎西洋的人，叫他带

了那难人一齐回本国去，又听见说那个长崎的人因有洋船，顺便托他寄信在球日本的人。我看起来琉球替日本原有交际往来，这事真真不错了。答：我们敝国原没有替他交际往来，那报知的事有些缘故。他飘来时节，那度

佳喇的人留在这里，他明明知道难人的事情，他回国之日把那洋人的事讲给日本的人知道，那日本的人传知洋人可就这样么。那书信到来的事也有些缘故，那度佳喇的人在敝国染病请医生调治，所以那长崎的朋友把个书信寄给病人问候，并没有替他结交往来。问：我听见说那英国的人来到东洋琉球两国已经啰唆，故此两国官民定要报

仇，同到哦啰嘶国学习鎗炮武艺。这个事情浙江官府奏明皇上，皇上叫他细细访问，真真这样么？答：我们小邦僻居海隅敌他不过的。那佛英的人留在敝国不能叫他回国，专赖

天朝的洪福，叫他无事回去，国家安宁，并没有学武报仇的事情。

\vfill
\begin{BookNotes}
\BookNote{7.}{度佳喇岛是位于南西诸岛北边吐噶喇列岛的总称，即以飞石状连结之所谓的七岛一口之岛、中之岛、诹访之濑岛、恶石岛、卧蛇岛、平岛及宝岛。1719年，尚敬受封之时，萨摩藩惟恐清朝以琉球受控于萨摩藩为由，断绝中琉间的宗藩关系，于是利用位于南西诸岛北部之吐噶喇列岛，虚构「宝岛」之名，并谎称所谓的日琉关系即是琉球与日本属岛—宝岛（度佳喇岛）之间存在的边境贸易关系，此举揭开了对中国隐蔽政策之序幕。虚构度佳喇岛除了不至于令琉中关系恶化以外，也是避免日中两国在琉球问题上发生冲突之措施。隐蔽政策本是出于岛津氏发起之琉球统治策下的计谋，但由于王府也因担忧受控于萨摩之实一旦暴露，进贡体制及王国的存续将会受到影响。因此，在隐蔽政策的推行问题上，王府及萨摩岛津氏双方意见达成一致，并将其作为幕藩体质下的共同政治课题，在整个近世期，一旦有漂流事件发生或册封使来琉之时，双方合力严密地展开各项隐蔽工作。(参考纸屋敦之《大君外交と東アジア》吉川弘文馆，1977年）。正史之《中山世谱》中关于「宝岛（度佳喇岛）」之存在有以下的记述。我国土瘠产少，国用不足。故与朝鲜日本暹罗爪哇等国，尝行通交之礼，互相往来，以备国用。万历年间，王受兵警，出在萨州。时王言，吾事中朝，义当有终。日本深嘉其志，卒被纵回。自尔而后，朝鲜日本暹罗爪哇等国互不相通。本国孤立，国用复缺。幸有日本属岛，度佳喇商民，至国贸易，往来不绝。本国亦得赖度佳喇，以备国用，而国复安然。故国人称度佳喇，曰宝岛。（《琉球史料丛书四》，东京美术，1972年，111页)。}
\end{BookNotes}

% 原书第 124 页（PDF 第 137 页）
\BookPage{124}{137}
\noindent 问：我听见讲你们琉球写个票文要买火轮船是真的么？8答：那火轮船价钱太贵，我们敝国没有金片银子，那火轮船其实买不得，况且

那火轮船替我敝国船只大不相同，那杠棋蓬舵不能使用，又没有煤炭，都是不便得。并没有要买的事。问：你们琉球要买那个船只有文书凭据，那里这样说。答：这个文书他国的人假借球人的名字写的，不是球人写的。

问：我听见讲琉球所管的外岛一个大岛、一个德岛、一个奇界岛、一个与论

岛、一个永良部岛，这个五岛昔年琉球屡屡遇着饥馑，就到那度佳喇岛借了好多大米，未曾还他。现今把那五岛的大米交纳那度佳喇岛，又有度佳喇的人设馆留在那里，这话是真的么。9答：这样事情不是真的。只要说明敝国穷苦的缘由就是替番洋的人这样说知。

虽说那度佳喇的人留在五岛，这也不是真的，我心里想起来那度佳喇的客商要做买卖，屡次留在那里，所以那个洋人这样说的。问：这个话虽是有道理，我听见说从此以前番洋的人在那外岛时候度佳喇的人

带了刀鎗铁砲阻隔，番洋走过的路上看到这个举动。那个五岛就是度佳喇所管的不错不错。答：不是所管的。那人阻隔的事情想必替番洋的人也有不和的事，所以这样

的。问：我听见说西洋人要做结好通商的事，球官辞说若是替他和好买卖，那度佳

喇的人往来绝迹，或者年成不好的时候借不得大米，都是饿死了。这话是

\vfill
\begin{BookNotes}
\BookNote{8.}{1851年岛津齐彬即藩王之位后，积极推行富国强兵政策，1857年10月砲术家市来四郎来琉，谎称是度佳喇岛的医师伊知良亲云上，经通事牧志朝忠介绍，翌年2月起开始与滞留琉球之法人传教士西拉鲁进行购买蒸汽军船的交涉，历时五个月，终于7月交涉达成协议，购买船身长60尺或4·50尺的蒸汽军船、大砲20门、新式步枪及其他装备，共价18万5000两，分六年付清。然而，由于齐彬的骤逝事态急变，此案全遭驳回。齐彬之死在琉球也掀起了一场政治斗争，即围绕三司官人事案引发的牧志·恩河事件。（上原兼善，《鎖国と藩貿易》，八重岳书房，1981年，299-300页)}
\BookNote{9.}{奄美诸岛被称为「道之岛，萨摩岛津氏侵琉之后被纳为萨摩的直辖地，但其表面上奄美诸岛的粮食收成所应缴纳的赋役仍归琉球所有，在奉幕府之命绘制《琉球国绘图》之时，亦将奄美诸岛归为琉球国领内。}
\end{BookNotes}

% 原书第 125 页（PDF 第 138 页）
\BookPage{125}{138}
真的么？答：不是真的。只要推辞结好通商的事，讲这个将就的话。问：我听见讲萨州前滨地方琉球船只多有湾泊时节，番洋兵船到来放了大砲烧

掉球船，你们把这些事情实实说给我知道。答：敝国船只没有烧掉。想必敝国外岛的船只飘到那里被他放砲烧了，这也定

不得的。近来又有番洋的船来到敝国，也不讲这个事。问：我听见说英国总兵亚国提督都到王殿相见球官，真真这样么？答:有这些事情。英官要到王殿相见总理官布政官，递给他的军机大臣的文

书。那亚官也要到王殿面谢备办他们物件的恩，又要庆贺新年的洪福，我们总理官等再三推辞说要在城外公馆相见，不要在王殿相见，那夷官人等都是不肯听从，直到王殿。那事情完了就是回去的。问：那文书的意思的怎么样呢？答:看那文书说是琉球接待伯德令的事须要尽行宾主的礼数，若有轻慢他，英

国要发义怒你们琉球，凡事任凭他的意，不要执拗。

问：我听见说那佛兰西亚米理干荷兰三国的人替你琉球议定条约打印给他，这

个事是真的么？10答：这些事情我们敝国再三推辞，他都是不肯听从，我们敝国没奈何议定的。

那亚荷两国的人或者他买了日用物件柴水，或者他的船只漂风损坏时节就替他修葺或者随便解送，不过这三样的事打印给他约定了。那佛国的人除了那个三件的事，又有一件事，他要借了地方盖起房屋随意住在。这个事情敝国严禁所以再三推辞，他们不肯听从，又推辞说我们敝国若有大事，不敢自专，奏闻皇上然后做得来的，故此要宽了年月，也是不肯，没奈何打印给他。从此以后，那国的人盖起房屋有久留的模样，敝国灾难不小，不能叫他回国，专赖天朝德威使他们回去了，此恩此德不能可报。

\vfill
\begin{BookNotes}
\BookNote{10.}{王府于1854年（咸丰四年）和美国签订琉美修好条约，翌年和法国签订琉法修好条约，1859年（咸丰九年）又与荷兰签订了琉兰修好条约。}
\end{BookNotes}

% 原书第 126 页（PDF 第 139 页）
\BookPage{126}{139}
\noindent 问：那个文书拿来给我看。答：这个事情我们回去转报管官。问：那三国夷人约定以后，来到这里时候应用物件自己买了么？也是球官料理

卖给他么？且这里不用金片银子，那卖给物价受了洋银么？答：我们敝国因有条约的事，许他自己到市上买了物件。若是市上没有物件，

管官替他料理卖给他们，我们敝国不用金片银子，所以把他洋银换了我的铜钱，他自己买了。问：这里金片银子不通买卖，那个洋银怎么样用得去呢？答：那个洋银都是藏在库里，没有费用，那伯德令回国时节分送他们。那佛夷

回国时候买了他住的房子，又兼屡屡垫补寄闽银子，其余银两今番要补垫官员货物的价钱。问：那洋银有多少呢？答：我听见管官说大约有一千二百元。问：那佛人借了地方做的房子发给价银么？答：那价银推辞不收。那佛人说若是不收价银我们不得做房子，须要受了价

银。没奈何把那人工木料的价钱收了洋银八百元，不收地方的价钱。问：那佛人回国时节买他住的房子多少价钱呢？答：把那洋银八百元就要买他房子，替他商量。他说我们住在好久，若是收了

原价没有道理，这也不肯的，所以把那洋银六百元买了他房。问：我听见说原来的亚船水梢做了强奸的事，这个事情细细说给我知道。答：那水梢经过那霸地方，他看人家关门，就是逾垣破开窗户强奸女人的，那

邻家的人并走路人等听见这事，走到那家赶出那水梢。那水梢慌忙逃走，众人忍耐不得丢丸赶打，那水梢魂飞魄散落了海中淹死。那强奸的人真真可恶，虽是立刻打他死了也没有得罪，不想那船提督说，审出犯人照法加罪。故此敝国狱官审问加罪。那提督又说，那水梢若是不死，罪罚不小，因他死了，这也没奈何的。那家女子给他强奸以后，心里十分怕羞，不吃


% 原书第 127 页（PDF 第 140 页）
\BookPage{127}{140}
饭死了。11问：我听说琉球通宝12的铜钱在东洋使用，这是真的么？答：这个铜钱我未曾看得。那外国假借琉球的名字新铸铜钱甚是可恶。我看

大清会典则礼13等书，私下铸钱天下的严禁，又兼外国另记年号也是违禁的，我们敝国不敢故犯大禁的。问：我听见西洋的人说你们琉球服从日本是真的么？答：不是这样。我们敝国地方褊小物件不多，原来替那日本属岛度佳喇人结交

通商买办进贡物件，又是买得日用物件，那度佳喇得人在日本收买那些东西卖给敝国，想必那西洋的人看得这个举动就说琉球在那日本的所管。我们敝国原来天朝的藩国，世世荷蒙封王，此恩此德讲不尽的，那有忘恩负义从服日本的道理。问：听见说琉球遇着江府将军即位的时节差官来朝见他们。那球官粧扮狠好，

带有牌板吹鼓，也有仗仪的绘图。我看起来琉球从服日本这也没有可疑。你们把这些事情实实说来。答：我们敝国在明朝时候，不论日本，凡是外国都有结交买卖。大清以后替外

\vfill
\begin{BookNotes}
\BookNote{11.}{此事件发生于贝理提督赴日本远征之时。贝理向王府要求对杀害美国人的犯人进行判处，对此，王府以虚构人物之「家满渡庆次」为犯人，判其终身流放八重山岛。负责处理此事件之官吏亦是虚设官职之布政官及琉球国中山府知府。(参照琉球大学附属图书馆藏之《八重山嶋江一世流刑手形》）}
\BookNote{12.}{琉球通宝的铸造是萨摩藩为重整财政所实施的政策。铸造的计划是由岛津齐彬于安政四年（咸丰七年、1857年）6月开始秘密进行，他一方面在御花园制炼所传授铸造法，另一方面为得到幕府的许可千方百计各方奔走。文久二年（同治元年、1862年）秋终于获准，但前提条件是二至三年内铸造量以一百万两为限，并须将其中一部分上缴幕府。同年12月24日在矶海滨建造大型的琉球通宝铸造所，开始铸造铜钱。萨摩藩意图混以天保钱交纳，于是仿照天保通宝钱，铸造了椭圆形的百文钱（1枚相当于124文），后又铸造圆形半朱钱（1枚相当于248文)，不久这些钱也流通于其他藩地。按从铸造到流通这一过程所需的各种费用以及周转资金的收益估算下来，琉球通宝实际铸造单价36文在流通时却相当于124-248文不等，获利莫大。因此，齐彬在位时原为铸造大砲而收集的佛具梵钟自不待言、连古制铜砲也均被挪为铸钱之用。据说梵钟之内甚至还有收自佐原城，刻有「日隅萨摩三州太守三位伊藤修理大夫」铭文之物。铜原料的匮乏，加之幕府末期海防、军备之需，铜价暴涨。受其影响，琉球社会原有的1比1的铜、铁钱兑换比率发生了大幅度的变化。而实际上，冠有琉球之名的琉球通宝并未在琉球国内流通。(参考东恩纳宽惇，《南島通貨志の研究》，《东恩纳宽惇全集4》第一书房，1979年，56-57页。参考《冲绳大百科事典》，冲縄夕イムス，1983年)}
\BookNote{13.}{「则礼」应是「则例」之误。}
\end{BookNotes}

% 原书第 128 页（PDF 第 141 页）
\BookPage{128}{141}
国不通，只有进贡天朝，多蒙天恩准了买来绸缎药料等项，国家使用没有不足，那敢朝见日本的官员。我们听见说往年大岛年成不好，到那度佳喇岛要借大米，那岛也是饥馑没有大米，那官员到江府替他商量借了许多大米借给大岛的人，那岛人民才得活命，所以那岛官员岛江府道谢。我们想起来外国的人把这个事情错认敝国差官来朝见的，我们敝国从古以来多蒙天朝的恩德，那有忘恩负义从服日本的道理。问：听见说琉球的人岛日本奇观的所在或者作诗或者做歌记载那国梓板，我看

起来琉球差官朝见的事也没有可疑。14答：敝国的人并没有作诗做歌的事。我们听见说往年大岛年成不好到度佳喇岛

要借大米，那里也是大饥没有大米。那里官员到了江府替他商量借了大米，转借大岛，所以到江府谢恩。想必那官员谢恩的时节做得诗歌纪名琉球的人，这也未可定。又有度佳喇的人屡屡带来日本有名的山川画儿，央托敝国的人看画做诗。敝国的人替他做诗写给他们，敢莫这样诗句记载梓板？且康熙年间靖南王谋叛时节，遣发游击陈应昌再三劝说丢弃大清从服靖王，我们国王不敢行那违天逆道的事，皇上晓得这个事赏给敕谕，褒嘉忠顺的。大人看见这事自然晓得不服日本的事。问：琉球原来穷苦，册封费用备办辛苦么？答：荷蒙皇上钦差封王的大典，又蒙两位大人的恩德，万物自然丰登，不到辛

苦。这是上下臣民原来所望的。问：遐迩贯珍云，琉球下等人家婚姻的时节不用媒人、不告父母、不行聘礼，

自己择了吉日做得夫妻，真真这样么？答：不是这样的。婚姻的礼关系最大，大约从依圣贤的法，只有富贵贫贱的

分。故此聘礼粧盒也有厚薄不等的，虽是下等人家也不敢做了那样无礼的事。

\vfill
\begin{BookNotes}
\BookNote{14.}{德川将军袭位或琉球国王袭封之时，琉球均须遣使前赴江户谒见将军，以表庆贺或谢恩。使节一行由那霸出港航抵萨摩山川港后，换船北上，经天草、呼户等地，由下关驶入濑户内海，中途停靠津轮地岛和明石港等地，抵达大坂后由淀川北上，于伏见登陆经由东海道前赴江户。赴江户之使节一行途经主要城市时，会特意以唐装打扮，并吹奏路次乐，途中亦与文人雅士吟诗和咏。王府派遣使节前赴江户之举始于1634年庆贺德川家光袭位之时，止于1850年尚泰即位之谢恩，共计18次。}
\end{BookNotes}

% 原书第 129 页（PDF 第 142 页）
\BookPage{129}{142}
\noindent 问：听见说琉球一个铜钱原来算用一文，如今怎么样扣作六文呢？15答：那个铜钱度佳喇的人在日本带来的，近来日本的铜钱很少使用不彀所以扣

作六文，那岛的人也是这样使用的。我们敝国若不是改用，他们不肯买卖，没奈何从他使用的。问：东西的价钱或者贵的或者贱的，都是因了年成好歹，这个不独琉球，虽是

他邦也是这样的，如今琉球的物价比前不同，都是太贵，是怎么样的？答：我们敝国铜钱扣作六文以后，物价也是太贵的，我们官府必要公平价钱争

奈徒然不贱的。从今以后渐渐低些的。

右者勅使様御寻之节御返答之条官话组仕候事

同治五年丙寅六月

\vfill
\begin{BookNotes}
\BookNote{15.}{为铸造大砲扩充军备以应对外国列强的压力，及为铸造琉球通宝以重整财政等原因，萨摩藩急需大量的铜材。结果造成了铜材的严重不足，从而导致了铜钱价值的大幅度上涨，铜钱和铁钱的兑换汇率发生了变化。虽然铜铁钱兑换汇率变动的指令是由琉球王府颁布，但事实上是配合了萨摩政策的实施。琉球王府于文久元年（咸丰十一年，1861）1月26日下达指令，将原本1比1之铜铁钱汇率改为铜钱1对铁钱2的比率。文久三年（同治二年，1863）3月10日改为1比4，元治元年（同治三年，1864）3月24日恢复为原来的1比2，元治元年（同治三年，1864）9月18日又变为1比3，庆应元年（同治四年，1865）2月2日升为1比4，庆应元年（同治四年，1865）6月28日竟高达1比6。4年半左右的时间内王府下达的汇率变动指令竟达6次之多，使得物价暴涨。(西里喜行，《清末中琉日關係史の研究》，京都大学出版会，2005年，248-249页。)}
\end{BookNotes}

% 原书第 130 页（PDF 第 143 页）
\BookPage{130}{143}
